\section{Background and Related Work}
\label{sec:background}

We group prior work by the question each body of work answers, because the
groupings are not competitors --- they occupy different positions in the same
stack, and \sys depends on several of them.

\subsection{ML lifecycle platforms}

The ML platform literature established that the model is not the deliverable: the
pipeline around it is. TFX~\cite{baylor2017tfx} articulated production-scale
lifecycle management; MLflow~\cite{zaharia2018mlflow} made experiment tracking,
model versioning, and artifact storage into general infrastructure; and
Kubeflow~\cite{kubeflow}, Airflow~\cite{airflow}, and DVC~\cite{dvc} covered
orchestration and data versioning. Work on automatic capture of experiment metadata and
provenance~\cite{schelter2017provenance} established that lineage has to be
recorded by the platform rather than by the practitioner --- a premise \sys inherits
and extends from experiments to releases. Alongside them, a diagnostic literature
described what goes wrong: hidden technical debt in ML
systems~\cite{sculley2015debt}, the engineering practices that
follow~\cite{amershi2019se4ml}, production-readiness rubrics~\cite{
breck2017mltest}, and documentation norms for models and
datasets~\cite{mitchell2019modelcards,gebru2021datasheets}.

\sys sits in this tradition and deliberately does not re-enter it. \mlflow{}'s
experiments, traces, assessments, prompt registry, artifact store, dataset
registry, and evaluation entry points are dependencies (D1,
\tabref{tab:decisions}). The distinction we draw is one of \emph{unit}: these
platforms govern models and pipelines, whose deployable is a serialized artifact
plus code. The artifacts that determine an LLM agent's behaviour are prompts, tool
definitions, skills, retrieval corpora, and judges --- and these have a property
models do not, which is that the deployable is often a \emph{selection} rather
than a binary. The consequence is that the interesting release primitive is alias
indirection, not artifact shipping, and none of the model-centric platforms
provide it for these types.

\subsection{Agent architectures and composition frameworks}

A large body of work establishes how to build agents:
ReAct~\cite{yao2023react} interleaving reasoning and acting,
Toolformer~\cite{schick2023toolformer} on tool invocation,
chain-of-thought prompting~\cite{wei2022cot},
self-critique and reflection~\cite{madaan2023selfrefine,shinn2023reflexion},
retrieval augmentation~\cite{lewis2020rag}, multi-agent
conversation~\cite{wu2024autogen}, and simulated societies of
agents~\cite{park2023generative}. The engineering counterpart is a set of
composition frameworks --- LangChain and LangGraph~\cite{langchain},
LlamaIndex~\cite{llamaindex}, CrewAI~\cite{crewai}, the OpenAI Agents
SDK~\cite{agentssdk} --- plus the Model Context Protocol~\cite{mcp2024} as an
emerging interoperability layer for tools and resources.

These answer \emph{how to build}, and \sys is explicitly not a competitor: an
agent composed with any of them is a client of a \sys-governed asset, resolving a
\livetarget{} at call time (\figref{fig:system}). The relationship is the one
between a web framework and a deployment system. \sys does interpret workflows
internally over a typed IR with \statNodeTypes{} node types, but that IR exists to
make a workflow a \emph{governable} artifact --- versioned, replayable, gated ---
rather than to compete on expressiveness.

\subsection{Prompt and program optimization}

Automated prompt search has moved quickly. APE~\cite{zhou2023ape} framed
instruction generation as search; OPRO~\cite{yang2024opro} used the model itself as
the optimizer; DSPy~\cite{khattab2024dspy} recast LLM pipelines as compilable
programs with MIPRO~\cite{opsahlong2024mipro} optimizing instructions and
demonstrations jointly; GEPA~\cite{agrawal2025gepa} showed reflective prompt
evolution competitive with reinforcement learning; and
TextGrad~\cite{yuksekgonul2024textgrad} backpropagates textual feedback.

\sys treats these as pluggable strategies behind one seam, and exposes
\statOptimizersW{} implemented paths (\figref{fig:pipeline}). The architectural
observation we would emphasize is that this literature has an implicit
precondition it does not supply: an optimizer needs a metric, a corpus, and a
baseline. Where those come from, whether they are versioned, and whether the
resulting candidate may be released are outside its scope --- and are exactly what
\secref{sec:capabilities} is about. This is why we characterize optimizers as
compilers rather than as lifecycle solutions.

\subsection{Evaluation, judges, and human alignment}

Holistic benchmark suites~\cite{liang2023helm} established multi-dimensional
evaluation as the norm; LLM-as-judge methods made per-example scoring
affordable~\cite{zheng2023judging,liu2023geval}; RAGAS~\cite{es2024ragas}
specialized this for retrieval-augmented systems; and harnesses such as
promptfoo~\cite{promptfoo} and OpenAI Evals~\cite{openaievals} made offline suites
routine; and Guardrails~\cite{guardrails} moved a subset of these checks inline, as
runtime validators rather than offline scorers --- a complementary position, since a
runtime validator constrains one call while an offline suite characterizes a
version. The known hazards are equally well established --- judges are biased
toward verbosity and toward outputs from related models, and agreement with human
labels must be measured rather than assumed, which is why \sys records
inter-rater agreement statistics for authored
judges~\cite{cohen1960kappa,fleiss1971kappa}.

The gap \sys addresses is not scoring quality but \emph{binding}. In the harness
model, a score is a number in a report. In \sys, a score is a durable record
attached to a candidate version, admissible as evidence at a gate, and carried
into a human review (\figref{fig:chain}). The same computation becomes a different
kind of object when it is bound to the artifact it describes.

\subsection{Observability and prompt-lifecycle platforms}

LangSmith~\cite{langsmith}, Langfuse~\cite{langfuse}, and Arize
Phoenix~\cite{phoenix} collect traces, run offline and online evaluations, and
present the results, converging on OpenTelemetry and its GenAI semantic
conventions~\cite{opentelemetry,semanticconv}. They are the closest neighbours to
\sys, so it matters to state the boundary accurately rather than sharply.

Two of them, together with \mlflow{}'s Prompt Registry, have grown past
observability into prompt lifecycle management, and a reader should not carry the
older mental model into this paper. As of 2026-08-03 all three provide versioned
prompts, a named pointer that client code resolves at call time --- \mlflow{}
aliases~\cite{mlflowpromptregistry}, LangSmith commit tags including a reserved
\code{production}~\cite{langsmithprompts}, Langfuse
labels~\cite{langfuseprompts} --- rollback by re-pointing, and role restrictions
on who may move the pointer. \mlflow{} also ships prompt optimization driven by a
dataset and scorers~\cite{mlflowoptimize}. \sys's mechanisms for prompts are
substantially the same mechanisms, and this paper does not claim otherwise.

What we did not find documented in any of them is a second governed artifact type
beyond prompts and models, or an evaluation result acting as a \emph{precondition}
on the pointer move rather than as an annotation beside it. Those two gaps, not an
empty ecosystem, are what \secref{sec:compare} argues \sys occupies, and
\tabref{tab:platforms} states the comparison row by row against dated sources.

\subsection{Isolation, policy, and provenance}

\sys executes operator-authored tool code and connects to third-party MCP servers,
which places it in a well-mapped security landscape. The foundational principles
are Saltzer and Schroeder's least privilege and fail-safe
defaults~\cite{saltzer1975protection}; the authority model is role-based access
control~\cite{sandhu1996rbac} with the integrity and separation-of-duty concerns
Clark and Wilson formalized~\cite{clark1987integrity}. Strong isolation for
untrusted code is a solved problem given the right substrate --- microVMs such as
Firecracker~\cite{agache2020firecracker}, or unprivileged sandboxing via
bubblewrap~\cite{bubblewrap} --- and policy engines such as Open Policy
Agent~\cite{opa} externalize authorization decisions. On the provenance side,
in-toto~\cite{torresarias2019intoto} and SLSA~\cite{slsa} define what it means for
an artifact's history to be attestable.

Two points of intellectual honesty follow, and we return to both. First, \sys's
local tool and MCP execution is \emph{containment} --- allowlists, a sanitized
environment, a private working directory, process timeouts --- and not isolation
in the Firecracker sense; we use the two words distinctly throughout
(\figref{fig:trust}). Second, \sys's provenance anchors are internal records, not
cryptographic attestations in the in-toto sense, and closing that gap is future
work (\secref{sec:future}). Prompt injection through retrieved or tool-returned
content~\cite{greshake2023injection,perez2022ignore} is a threat \sys's governance
model constrains --- by bounding what a capability may do --- but does not solve.

\subsection{Distributed-systems foundations}

\sys's execution model rests on old results. Its durable queues are status,
claim, and lease columns on relational rows, which is the transactional-outbox and
lease pattern rather than anything new~\cite{gray1981transaction,kleppmann2017ddia}.
Its two cross-store crossings cannot be made atomic, so they are compensated
rather than transactional, in the spirit of sagas~\cite{garciamolina1987sagas} and
of Helland's argument that distributed transactions are the wrong default for
scalable systems~\cite{helland2007apostates}. Its SLO and incident surface follows
standard site-reliability practice~\cite{beyer2016sre}. We claim no novelty in any
of this; we claim only that applying it to \emph{agent-resource release} is what
turns an observability tool into a control plane, and \secref{sec:design} records
where each borrowed pattern charges its price.
